\documentclass[10pt,a4paper]{amsart}
\usepackage[margin=19mm]{geometry}
\usepackage{amsmath,amssymb,mathtools}
\usepackage[scr=rsfs,cal=euler]{mathalfa}
\usepackage{xfrac}
\usepackage[unicode,hidelinks,bookmarks=false]{hyperref}
\newcommand{\ie}{i.e.\ }
\newcommand{\cf}{cf.\ }
\newcommand{\resp}{respectively\ }
\newcommand{\Subsection}{\subsection}
\providecommand{\nobreakdash}{\nobreak\hbox{-}}
\setlength{\emergencystretch}{2em}
\multlinegap0pt
\usepackage{bm}
\usepackage{bbm}
\usepackage{dsfont}
\usepackage{xspace}
\usepackage{cancel}
\usepackage{mathtools}
\usepackage{amsthm}
\makeatletter
\@ifundefined{thm}{\newtheorem{thm}{Theorem}}{}
\makeatother
\title[Machine learning moment closure models for the radiative tra]{Machine learning moment closure models for the radiative transfer equation IV: enforcing symmetrizable hyperbolicity in two dimensions}
\author{Juntao Huang}
\date{Source: arXiv:2604.20143v1 (CC BY 4.0)}
\begin{document}
\maketitle

\noindent\emph{Source and licence.} Machine learning moment closure models for the radiative transfer equation IV: enforcing symmetrizable hyperbolicity in two dimensions. Version arXiv:2604.20143v1 (CC BY 4.0), distributed under the Creative Commons Attribution 4.0 International licence, as recorded in the arXiv licence metadata for this exact version and shown on the abstract page https://arxiv.org/abs/2604.20143v1. Full rights notice: licenses/arxiv-2604.20143-CC-BY-4.0.txt.\par\smallskip

\noindent\textit{Editorial scope.} Counted unit: the Thm whose source label is \texttt{thm:real-sph-harm-orthonormal} in the article (source0.tex lines 871--895, source label \texttt{thm:real-sph-harm-orthonormal}), delivered together with the article's own complete original proof of it (source0.tex lines 896--951, one \texttt{proof} environment of 56 lines). No mathematics is written, repaired or re-derived: statement and proof are the article's own text line for line. Required context retained uncounted: none. Every definition, display and label of those lines is retained unchanged. Omitted from this package and not redistributed: 1--870, 952--1253. Nothing inside a retained excerpt is omitted or altered: every retained body is delivered exactly as the article's lines are written, so no omission receipt of any kind accompanies this package. Citations: the retained excerpts carry no citation command at all, so no bibliography entry of the article is transcribed and no reference is redistributed. Reference adaptation: none was needed and none was made; every reference inside the delivered unit resolves inside it, so no unresolved reference or citation is produced. Printed slips kept literal: no printed word, symbol, hypothesis or number of the counted statement or of its proof was altered. Layout and class: the article's own class is \texttt{article} with packages including amsmath, amssymb, amsthm, amsfonts, bm, geometry, graphicx, hyperref, subcaption; the wrapper is the lane's pinned \texttt{amsart} editorial wrapper at 10pt on A4, which restates the theorem-like environments the retained text needs and adds the article's own macro definitions that the retained text needs (none), with extra wrapper packages (bm, bbm, dsfont, xspace, cancel, mathtools). The delivered numbering is the unit's own. Assets: no retained span contains a figure environment, a graphics-inclusion command, a PDF-page inclusion, a table or a TikZ picture; no image file is copied into this package, the package declares no asset dependency and no image file is redistributed with it. Licence: version arXiv:2604.20143v1 of this article is distributed under the Creative Commons Attribution 4.0 International licence, as recorded in the arXiv licence metadata for this exact version and shown on the abstract page https://arxiv.org/abs/2604.20143v1. A full rights notice is delivered at licenses/arxiv-2604.20143-CC-BY-4.0.txt. Characters: every retained excerpt is pure ASCII, so the delivered engine, XeLaTeX with its default Latin Modern text font, reports no missing character.
\begin{thm}[Real spherical harmonics form an orthonormal basis]\label{thm:real-sph-harm-orthonormal}
Let
\[
N_{l,m}:=\sqrt{\frac{2l+1}{4\pi}\frac{(l-m)!}{(l+m)!}},
\qquad l\ge 0,\quad 0\le m\le l,
\]
and define
\[
\Phi_l^0(\mu,\varphi):=N_{l,0}P_l(\mu),
\]
and, for \(1\le m\le l\),
\[
\Phi_l^{m,\mathrm c}(\mu,\varphi)
:=\sqrt{2}\,N_{l,m}P_l^m(\mu)\cos(m\varphi),
\qquad
\Phi_l^{m,\mathrm s}(\mu,\varphi)
:=\sqrt{2}\,N_{l,m}P_l^m(\mu)\sin(m\varphi).
\]
Then the family
\[
\bigl\{\Phi_l^0,\ \Phi_l^{m,\mathrm c},\ \Phi_l^{m,\mathrm s}
:\ l\ge 0,\ 1\le m\le l\bigr\}
\]
is an orthonormal basis of \(L^2(\mathbb S^2)\).
\end{thm}

\begin{proof}
    We first derive the integration formula for the surface integral on the sphere with the variables \((\mu,\varphi)\). With the spherical coordinates \((\theta,\varphi)\) with $\Omega = (\sin\theta\cos\varphi,\sin\theta\sin\varphi,\cos\theta)$, the surface integral of a function \(f:\mathbb S^2\to\mathbb R\) can be written as
    \begin{equation}
    \int_{\mathbb S^2} f(\Omega)\,d\Omega = \int_0^{2\pi}\int_0^\pi f(\theta,\varphi)\sin\theta\,d\theta\,d\varphi.
    \end{equation}
    Now with the change of variable $\mu=\cos\theta$, we have $d\mu=-\sin\theta\,d\theta$, so that
    \[
    \int_{\mathbb S^2} f(\Omega)\,d\Omega
    =
    \int_0^{2\pi}\int_0^\pi f(\theta,\varphi)\sin\theta\,d\theta\,d\varphi
    =
    \int_0^{2\pi}\int_{-1}^1 f(\mu,\varphi)\,d\mu\,d\varphi.
    \]

    We now check the orthonormality in the following cases: 
    \begin{enumerate}

    \item 
    The orthogonality of (\(\Phi_l^0\), \(\Phi_{l'}^{m,\mathrm c}\)), (\(\Phi_l^0\), \(\Phi_{l'}^{m,\mathrm s}\)), and (\(\Phi_{l'}^{m,\mathrm c}\), \(\Phi_{l'}^{m,\mathrm s}\)) can be easily checked by the orthogonality of the trigonometric functions of \(\varphi\).

    \item 
    The orthonormality of $\Phi_l^0$ and $\Phi_{l'}^0$ can be checked by directly computing
    \begin{equation}
    \begin{aligned}
        \int_{\mathbb S^2}\Phi_l^0\Phi_{l'}^0\,d\Omega ={}&
        2\pi N_{l,0}N_{l',0}\int_{-1}^1 P_l(\mu)P_{l'}(\mu)\,d\mu \\
        ={}& \frac{2\pi}{4\pi}\sqrt{(2l+1)(2l'+1)}\int_{-1}^1 P_l(\mu)P_{l'}(\mu)\,d\mu \\
        ={}& \frac{2\pi}{4\pi}\sqrt{(2l+1)(2l'+1)} \frac{2}{2l+1} \delta_{ll'} = \delta_{ll'}.
    \end{aligned}
    \end{equation}
    Here we use the orthogonality relation of the Legendre polynomials:
    \begin{equation}
        \int_{-1}^1 P_l(\mu)P_{l'}(\mu)\,d\mu = \frac{2}{2l+1}\delta_{ll'}.
    \end{equation}

    \item 
    The orthonormality of \(\Phi_l^{m,\mathrm c}\) and \(\Phi_{l'}^{m',\mathrm c}\) can be checked by directly computing
    \begin{equation}
    \begin{aligned}
        \int_{\mathbb S^2}\Phi_l^{m,\mathrm c}\Phi_{l'}^{m',\mathrm c}\,d\Omega ={}&
        2N_{l,m}N_{l',m'}\Bigl(\int_{-1}^1 P_l^m(\mu)P_{l'}^{m'}(\mu)\,d\mu\Bigr)\Bigl(\int_0^{2\pi}\cos(m\varphi)\cos(m'\varphi)\,d\varphi\Bigr) \\
        ={}& 2N_{l,m}N_{l',m'} \frac{2(l+m)!}{(2l+1)(l-m)!}\delta_{ll'} \pi \delta_{mm'} \\
        ={}& \delta_{ll'}\delta_{mm'}.
    \end{aligned}
    \end{equation}
    Here we use the orthogonality relation of the associated Legendre polynomials:
    \begin{equation}
        \int_{-1}^1 P_l^m(\mu)P_{l'}^m(\mu)\,d\mu = \frac{2(l+m)!}{(2l+1)(l-m)!}\delta_{ll'}
    \end{equation}
    Similar computation gives
    \begin{equation}
    \int_{\mathbb S^2}\Phi_l^{m,\mathrm s}\Phi_{l'}^{m',\mathrm s}\,d\Omega = \delta_{ll'}\delta_{mm'}.
    \end{equation}
    The proof is complete.
\end{enumerate}
\end{proof}

\end{document}
